\section{Diapositivas y evidencia visual}

\subsection{Capas de estrategia y evidencia: Slide 0-04}

\begin{figure}[H]
\centering
\includegraphics[width=0.8\textwidth]{slide-0-04.jpg}
\caption{Slide 0-04 superpone el reward model, la policy update y el proceso de gobernanza en un mismo lienzo, y enfatiza el track de evidence-to-action.}
\end{figure}

Slide 0-04 muestra cómo se recopilan y revisan varias evidence signal dentro de una misma timeline. Recuerda a los equipos que no pueden centrarse solo en un benchmark, sino que deben convertir la evidence matrix en un dashboard operativo.

\begin{knowledgebox}{Evidence dashboard por capas}
Si en el dashboard se presentan benchmark, human eval y deployment logs en una stacked view, es más fácil identificar en qué capa la signal hace drift primero.
\end{knowledgebox}

\subsection{Ritmo operativo: Slide 0-05}

\begin{figure}[H]
\centering
\includegraphics[width=0.8\textwidth]{slide-0-05.jpg}
\caption{Slide 0-05 traza el ritmo de alignment-engineering (data → training → deployment → monitoring) junto con los action items.}
\end{figure}

Cada etapa cuenta con sus artifacts: la data stage tiene un prompt catalog; la training stage, un KL margin report; la deployment stage, una drift alarm; y la monitoring stage, un ops playbook. Un diagrama de ritmo como este puede servir de referencia para la quarterly planning.

\begin{importantbox}{Lo esencial del ritmo operativo}
La vista timeline+artifact permite que los equipos multifuncionales lleguen rápidamente a un consenso en las reuniones de planning. Así se evita que la lógica de alignment se rompa entre la entrega y el monitoreo.
\end{importantbox}

\subsection{Matriz de calidad: Slide 0-06}

\begin{figure}[H]
\centering
\includegraphics[width=0.8\textwidth]{slide-0-06.jpg}
\caption{Slide 0-06 organiza en una matriz la calidad, la confianza y la gobernanza, y enfatiza los criterios de medición de cada stakeholder.}
\end{figure}

Esta diapositiva combina quality gate, trust signal y governance checkpoints en una matriz. Así es más fácil transmitir la evidence capa por capa a product, ops y legal.

\begin{warningbox}{No reforzar una signal de forma aislada}
Reforzar solo una automation metric (como la latency) sin integrarla con una governance checklist lleva a lanzar versiones eficientes pero inseguras. Slide 0-06 nos recuerda que debemos mantener el equilibrio dentro de la matriz.
\end{warningbox}

\subsection{Resumen de la sección}

Estas diapositivas sirven como visual anchor: convierten el RLHF pipeline abstracto en un evidence-dashboard layer-by-layer, un ops rhythm y una governance matrix, lo que favorece la cross-team alignment.
